\chapter{Wi-Fi Mood Lamp}

\section{Overview}
In this project, your ESP32-C3 becomes both a Wi-Fi access point and a light controller.
After uploading a Python file plus one HTML dashboard file, you will connect your phone or
laptop to the microcontroller's network, open a local web page, and control a chain of
three WS2812B LEDs.

You will build:
\begin{itemize}
    \item A password-protected Wi-Fi network hosted by the ESP32-C3
    \item A local web dashboard served directly from the board
    \item Manual lighting controls for all LEDs together and for each LED individually
    \item A set of built-in animations (breathing, rainbow, wipe, and blink)
\end{itemize}

By the end, you will have a complete wireless ``mood lamp'' that works without internet.

\pagebreak

\section{Components}

\begin{components}
    \item Seeed XIAO ESP32-C3 microcontroller
    \item 3-pixel WS2812B LED strip
    \item Breadboard
    \item Jumper wires
\end{components}

\section{Directions}

\subsection{Creating the circuit}
Before starting, remove components from earlier projects that are no longer needed.

% Seating the microcontroller. This is identical for every project, so the
% coordinates live here rather than in any one project chapter. The diagram
% generator reads the \bbhole definitions below for every breadboard diagram
% it draws, which is why this file is the only place they appear.
\subsubsection{Attach the microcontroller to the breadboard}
Orient the microcontroller so that the pin labeled \textbf{5V} goes into hole \bbhole{mcu.5v}{H1}
(Column H, Row 1), \textbf{GPIO2} into \bbhole{mcu.gpio2}{D1}, \textbf{GPIO20} into
\bbhole{mcu.gpio20}{H7}, and \textbf{GPIO21} into \bbhole{mcu.gpio21}{D7}. Then press it down---this
takes more force than you might expect.

\begin{figure}[H]
    \centering
    \maybeincludegraphics[width=.95\linewidth]{common/microcontroller_seated_in_breadboard.png}
    \caption{So far, so good!}
\end{figure}


\subsubsection{WS2812B strip wiring}
The LED strip has a right-angle header so it stands upright once seated.
Push those three pins into the top half of the board: \textbf{5V} into
\bbhole{pixels.header.power}{J12}, \textbf{DI} into \bbhole{pixels.header.data}{J13}, and
\textbf{GND} into \bbhole{pixels.header.ground}{J14}. Then add these connections:
\begin{center}
\begin{tabular}{|l|l|l|}
    \hline
    \textbf{WS2812B pin} & \textbf{XIAO pin} & \textbf{Jumper holes} \\
    \hline
    \connectionrow{pixels-power}{5V}{5V}{I1}{I12}
    \connectionrow{pixels-data}{DI}{GPIO 20}{I7}{I13}
    \connectionrow{pixels-ground}{GND}{GND}{I2}{I14}
    \hline
\end{tabular}
\end{center}

\begin{figure}[H]
    \centering
    \maybeincludegraphics[width=.95\linewidth]{project_5/wiring.png}
    \caption{Connect the three-pixel LED strip to the XIAO ESP32-C3.}
\end{figure}

\subsection{Programming the microcontroller}
Upload both files from the project code folder to your microcontroller:
\begin{itemize}
    \item \texttt{project\_5\_wifi\_mood\_lamp.py}
    \item \texttt{resources/project\_5\_wifi\_mood\_lamp.html}
\end{itemize}

Keep the \texttt{resources/} folder name exactly the same on the microcontroller filesystem.
The Python script loads the dashboard from that path.

Then open \texttt{project\_5\_wifi\_mood\_lamp.py} in the IDE and run it.
Refer back to Chapter \ref{ide} for IDE instructions.

When the script starts, watch the serial output. It prints:
\begin{itemize}
    \item A unique Wi-Fi network name (SSID), such as \texttt{MoodLamp-1A2B3C}
    \item The Wi-Fi password defined in code
    \item The local IP address for the control page
\end{itemize}

From your laptop or phone:
\begin{enumerate}
    \item Open Wi-Fi settings and connect to the printed SSID
    \item Enter the password
    \item Open a browser and navigate to the printed IP address (for example \texttt{http://192.168.4.1/})
\end{enumerate}

The dashboard includes:
\begin{itemize}
    \item \textbf{Global controls:} Set color and brightness for all LEDs at once
    \item \textbf{Per-LED controls:} Set each LED independently
    \item \textbf{Animation controls:} Start/stop breathing, rainbow, wipe, and blink
\end{itemize}

\begin{tcolorbox}[colback=green!5!white,colframe=green!40!black]
    \textbf{Tip:} If many students are in the room, the unique suffix in the SSID helps you identify
    your own board quickly.
\end{tcolorbox}

\section{Review}
In this project, you combined embedded networking and LED control in one program. You practiced:
\begin{itemize}
    \item Hosting a Wi-Fi access point directly on an ESP32-C3
    \item Serving a local web interface from a microcontroller
    \item Parsing user commands from HTTP requests
    \item Driving WS2812B LEDs with both static colors and animated effects
    \item Designing a system that is interactive, wireless, and real-time
\end{itemize}

\section{Possible Extensions}
\begin{itemize}
    \item Add a ``favorites'' list of saved color presets
    \item Add a timed auto-off feature (for example, sleep after 30 minutes)
    \item Add a ``party mode'' that randomly cycles between multiple effects
    \item Add an ambient light sensor and auto-adjust brightness
    \item Add a physical button that toggles between saved scenes without opening the web page
\end{itemize}
